\documentclass{article}
\usepackage[UTF8]{ctex}
\usepackage{amsmath,amssymb,amsfonts}
\usepackage{enumitem}
\usepackage{geometry}
\geometry{a4paper, margin=1in}
% === 错题本排版优化（不参与题目标识） ===
\setlength{\parindent}{0pt}
\setlength{\parskip}{0.35em plus 0.10em minus 0.08em}
\setlength{\emergencystretch}{3em}
\setlist{itemsep=0.25em, topsep=0.40em, parsep=0pt}
\allowdisplaybreaks[2]
% === 错题本排版优化结束 ===
\begin{document}
\section*{云南大学2024年高等代数真题}
\section*{183．云南大学2024年高等代数试卷}
一．填空题．每题5分，共30分．\\

\subsection{题目1}
1．多项式 \(\left(x^{2023}-1, x^{2024}-1\right)=\) \(\_\_\_\_\)．

\subsection{题目2}
2．秩为 \(r\) 的 \(n\) 阶方阵 \(A\) 满足 \(A^{2}=A\) ，则 \(A\) 的所有一阶主子式之和 = \(\_\_\_\_\)．

\subsection{题目3}
3．实二次型 \(x_{1}^{2}-3 x_{3}^{2}-2 x_{1} x_{2}+2 x_{1} x_{3}-2 x_{2} x_{3}\) 的符号差为 \(\_\_\_\_\)．

\subsection{题目4}
4．由所有 \(n\) 阶对称矩阵构成的实线性空间的维数 = \(\_\_\_\_\)．

\subsection{题目5}
5．若3阶方阵 \(A\) 的行列式是1，迹是2，并且多项式 \(f(x)=x^{2}-2 x\) 与 \(A\) 的特征多项式不互素，写出\(A\) 的所有特征值 \(\_\_\_\_\)．

\subsection{题目6}
6．记 \(E\) 是 \(n\) 阶单位阵，若 \(n\) 阶方阵 \(A, B\) 满足 \(A B-B A=E+A\) ，则 \(A^{2024}\) 的行列式等于 \(\_\_\_\_\)．

二．（15 分）求实数 \(b\) ，使得实系数线性方程组

\[
\left\{\begin{array}{l}
x_{1}+x_{2}+x_{3}+x_{4}=2 \\
2 x_{1}+2 x_{3}=0 \\
2 x_{1}+x_{2}+2 x_{3}+x_{4}=2 \\
b x_{1}+2 x_{2}+b x_{3}+2 x_{4}=b
\end{array}\right.
\]

有解，并在有解的情况下求所有的解。

三．（15 分）计算 \(n\) 阶矩阵 \(\left(\begin{array}{ccccc}1 & 2 & 3 & \cdots & n \\ 2 & 2 & 3 & \cdots & n \\ 3 & 3 & 3 & \cdots & n \\ \vdots & \vdots & \vdots & & \vdots \\ n & n & n & \cdots & n\end{array}\right)\) 的行列式。

四．（15 分）设矩阵 \(A=\left(\begin{array}{llll}2 & 0 & 2 & 4 \\ 0 & 2 & 0 & 2 \\ 0 & 0 & 2 & 0 \\ 0 & 0 & 0 & 2\end{array}\right)\) ．求 \(A\) 的 Jordan 标准形．

五．（15 分）用正交的线性替换化四元实二次型 \(\sum_{1 \leq i \leq j \leq 4} x_{i} x_{j}\) 成平方和的形式．

六．（15 分）证明：每一个 \(n\) 维线性空间都可以写成 \(n\) 个 1 维线性空间的直和．

七．（15 分）证明：\(n\) 阶实对称矩阵

\[
A=\left(\begin{array}{ccccc}
2 & -1 & & & \\
-1 & 2 & -1 & & \\
& -1 & 2 & \ddots & \\
& & \ddots & \ddots & -1 \\
& & & -1 & 2
\end{array}\right)
\]

正定，并且它有 \(n\) 个互不相同的正特征值．

八．（15 分）设 \(V\) 是 \(n\) 维欧氏空间，\(\gamma, \alpha_{1}, \cdots, \alpha_{m} \in V\) 满足 \(\left(\gamma, \alpha_{k}\right)>0,\left(\alpha_{i}, \alpha_{j}\right) \leq 0\) 对一切的\(i, j, k \in\{1,2, \cdots, m\}\) 且 \(i \neq j\) 。证明：\(\alpha_{1}, \cdots, \alpha_{m}\) 线性无关。

九．（15 分）记 \(\operatorname{End}(V)\) 是由 \(n\) 维复线性空间 \(V\) 上的所有线性变换构成的线性空间．固定 \(\sigma \in \operatorname{End}(V)\)，定义 \(\operatorname{End}(V)\) 上的变换 \(\varphi\) 为 \(\varphi(\rho)=\sigma \rho-\rho \sigma, \forall \rho \in \operatorname{End}(V)\) 。证明：\(\varphi\) 是 \(\operatorname{End}(V)\) 上的线性变换，并且若线性变换 \(\tau \in \operatorname{End}(V)\) 使得对某个正整数 \(m\) 有 \(\varphi^{m}(\tau)=0\) ，则 \(\sigma\) 的任一个根子空间都是 \(\tau\)－子空间。
\end{document}
